\scanpage{13}
\mspage{4}
Je considère maintenant la sous catégorie $ {\unsure{\mathscr B|\mathcal B|B}}_0$ de $\widehat M\supset M$, formée des objets $X$ de la forme suivante :
\[
\left(\begin{aligned}
&X\text{ est }\deleted{\text{\unsure{lim|limite}}}\inserted{\text{sup}}\text{ de ses sous-objets isomorphes \unsure{à|:} des}\\
&\deleted{\text{\unsure{ses|des}}}\qquad\text{objets de }M\text{ (modèles)}
\end{aligned}\right.
\]
% slanted
\margin{Cela \unsure{signifie|implique} que \unsure{tt} hom \deleted{de} $S\unsure{-|\to}X$, où $S\in\operatorname{Ob}M$, est un mono}

Si $K,L\subset X$ sont des sous-objets modèles, alors $K\cap L$ est sup des sous-objets de $K\cap L$ qui sont des modèles

Soit \unsure{donc|Donc} $\widetilde X$ l'ens des \deleted{\unsure{obj}} sous-objets de $X$ qui sont des modèles. Alors $\widetilde X$. est un ens ordonné \inserted{$K$}, muni de la structure suivante \inserted{(``$M$-épinglage'' de $K$)} : \deleted{\unsure{à|l'} \unsure{savoir|une}}

a) \deleted{\unsure{un|une} \unsure{de} \unsure{Pour}} Une application strictement croissante
\[
K\xrightarrow{\tau_K}\deleted{\unsure{I|N}}\unsure{N|\mathbb N}
\qquad\text{(application ``\uline{type}'')}
\]

b) Pour \unsure{tt} $\xi\in K$, \inserted{posant $n=\tau_K(\unsure{x|\xi})$,} \unsure{une|un} \deleted{\unsure{applic\textsuperscript{n}|applic\textsuperscript{t}}} \unsure{isomorphis\textsuperscript{me}|isomorphisme}
\[
\deleted{\unsure{D|I}}\;{\unsure{I|J}}_n\underset{u_\xi^K}{\simeq}K_{\leq\xi}
\qquad\text{(application d'épinglage de }K_{\leq\unsure{\xi|\eta}}\text{)}
\]
% slanted
\margin{les applications $I_n\unsure{-|\to}K$ définies par les $\xi\in K(n)=\unsure{(}\tau^{-1}(n)$ sont dite \uline{admissibles}}
\[
u_\xi^K(d_n)=\xi\quad\text{(= plus grd élément de }K_{\leq\xi}\text{)}
\]
avec \deleted{condition} les propriétés \deleted{de transitivité} :
% archive sheet number: 12
